%% toc: P15 174 174.9,174.10
\section{Proofs: assignment as a structural model}
\label{sec:annex-p15-w02b-p15-077-assignment-scm}

Three finite facts are used throughout. All random variables below are
functions on a finite outcome space \(\Omega\) with mass \(q\).
\begin{enumerate}
\item[(F1)] If \(V\) takes values in a finite set with law \(p\), then
      \(\mathbb E[F(V)]=\sum_vp(v)F(v)\) for every real \(F\): group
      the finite sum \(\sum_\omega q(\omega)F(V(\omega))\) by the value
      \(v=V(\omega)\), and note \(\sum_{\omega:V(\omega)=v}q(\omega)=p(v)\).
\item[(F2)] Suppose finite random variables \(A_1,\ldots,A_m\) satisfy
      \[
       \mathbb P(A_1=a_1,\ldots,A_m=a_m)=\prod_k\mathbb P(A_k=a_k)
      \]
      for all values. Then
      \(\mathbb E[\prod_k\alpha_k(A_k)]=\prod_k\mathbb E[\alpha_k(A_k)]\)
      for real functions \(\alpha_k\). By (F1) applied to the tuple, the
      left side is a finite multiple sum of
      \(\prod_k\alpha_k(a_k)\mathbb P(A_k=a_k)\), which factors into the
      product of the single sums.
\item[(F3)] If \(A\) is independent of a finite tuple \(T\) and
      \(W=\varphi(T)\), then \(A\) is independent of \(W\):
      \(\mathbb P(A=a,W=x)=\sum_{s:\varphi(s)=x}\mathbb P(A=a,T=s)
      =\mathbb P(A=a)\sum_{s:\varphi(s)=x}\mathbb P(T=s)
      =\mathbb P(A=a)\mathbb P(W=x)\).
\end{enumerate}
A constant random variable is independent of every random variable, since
the joint probability with value \(c\) equals the other marginal.

\begin{annexitem}{annex:p15-w02b-p15-077-corr}{thm:p15-w02b-p15-077-correspondence}{Structural form}
\begin{proof}
The graph is acyclic: every edge leaves vertex \(0\), and no edge enters
it, so no directed path returns to its start. The order
\(0,1,\ldots,N\) is topological.

\emph{Part 1.} At \(\omega\), vertex \(0\) has no parents, so
\(V_0(\omega)=f_0(U_0(\omega))=\omega\). Then
\(V_i(\omega)=f_i(V_0(\omega),0)=Y_i(\omega_i)\). By
\cref{thm:p16-causal-solution} this is the unique solution. For
\((z,y)\in S_0\times\prod_iS_i\), the outcomes with
\(V(\omega)=(z,y)\) must satisfy \(\omega=V_0(\omega)=z\). The single
candidate \(\omega=z\) qualifies exactly when \(y_i=Y_i(z_i)\) for every
\(i\). Hence \(p_{M_E}(z,y)=q(z)=\varpi(z)\) in that case and \(0\)
otherwise. In \(E\), on the event \(\{Z=z\}\) the recorded vector is
\((Y_i(z_i))_i\), so
\(\Pr(Z=z,(Y_i(Z_i))_i=y)\) is \(\varpi(z)\) when
\(y=(Y_i(z_i))_i\) and \(0\) otherwise. The two laws agree.

\emph{Part 2.} By \cref{def:p16-causal-intervention}, the model
\(M_{0=z}\) replaces only the equation of vertex \(0\), by the constant
\(z\), and keeps \(\Omega\), \(q\), the disturbances and every \(f_i\).
Evaluating in the same order gives \(V_0^{0=z}(\omega)=z\) and
\(V_i^{0=z}(\omega)=f_i(z,0)=Y_i(z_i)\), the unique solution by
\cref{thm:p16-causal-intervention-laws}. Taking \(z=t\mathbf 1\) gives
\(V_i^{0=t\mathbf 1}=Y_i(t)\).

\emph{Part 3.} By Part~2,
\(h(V_1^{0=t\mathbf 1},\ldots,V_N^{0=t\mathbf 1})\) is the constant
\(N^{-1}\sum_iY_i(t)\), and the expectation of a constant is that constant
because \(\sum_\omega q(\omega)=1\). Interchanging the finite sums,
\[
 \tau=\frac1N\sum_i\sum_tw_tY_i(t)
     =\sum_tw_t\frac1N\sum_iY_i(t)=\sum_tw_t\mu_{M_E}(t)=\tau(M_E).
\]
If \(\sum_tw_t=0\), subtracting
\(\bigl(\sum_tw_t\bigr)\mu_{M_E}(t_0)=0\) gives
\(\tau=\sum_tw_t(\mu_{M_E}(t)-\mu_{M_E}(t_0))\), in which the term
\(t=t_0\) vanishes. Each bracket is
\(\mathbb E[h(V^{0=t\mathbf 1}_{\{1,\ldots,N\}})]
 -\mathbb E[h(V^{0=t_0\mathbf 1}_{\{1,\ldots,N\}})]\), a difference of
intervention means of the real response function \(h\) at two fixed
interventions on \(I=\{0\}\). This is the average intervention effect of
\cref{def:p16-causal-intervention}. For the weights \(1\) at \(t\),
\(-1\) at \(s\) and \(0\) elsewhere, take \(t_0=s\); only the term
\(t\) remains, with coefficient one.

\emph{Part 4.} The same order evaluates \(M\): at each \(\omega\),
\(V_0(\omega)=U_0(\omega)\) and \(V_i(\omega)=g_i(U_0(\omega)_i)\).
The event \(\{V=(z,y)\}\) is therefore \(\{U_0=z\}\) when
\(y_i=g_i(z_i)\) for every \(i\), and empty otherwise. Its probability is
\(\varpi(z)\) or \(0\), respectively, which is the law of Part~1 for the
potential responses \(Y_i(t)=g_i(t)\). Under
\(\operatorname{do}(V_0=z)\), both models have the solution
\((z,(g_i(z_i))_i)\) at every outcome, so both intervention laws are point
mass at this tuple. Mass functions are compared on
\(\mathcal T^N\times\mathbb R^N\), taking value zero outside the value sets.
\end{proof}
\end{annexitem}

\begin{annexitem}{annex:p15-w02b-p15-077-ident}{thm:p15-w02b-p15-077-identification}{Identifying functional}
\begin{proof}
Let \(M\in\mathcal C_{\mathrm{NI}}(\varpi)\). The graph and order are as
in the preceding proof. Write \(U_{0,i}\) for the \(i\)th coordinate of
\(U_0\).

\emph{Part 1.} Fix \(t\) and \(i\). Under
\(\operatorname{do}(V_0=t\mathbf 1)\) the equation of vertex \(i\) is
unchanged, so \(V_i^{0=t\mathbf 1}=f_i(t\mathbf 1,U_i)=g_i(t,U_i)\).
By linearity of finite expectation,
\(\mu_M(t)=N^{-1}\sum_i\mathbb E[g_i(t,U_i)]\).
Observationally \(V_0=U_0\) and \(V_i=g_i(U_{0,i},U_i)\). By (F1)
with \(F(z,y)=\mathbf 1_{\{z_i=t\}}y_i\),
\[
 \sum_{(z,y)}p_M(z,y)\mathbf 1_{\{z_i=t\}}y_i
 =\mathbb E\bigl[\mathbf 1_{\{U_{0,i}=t\}}g_i(U_{0,i},U_i)\bigr]
 =\mathbb E\bigl[\alpha(U_0)\beta(U_i)\bigr],
\]
where \(\alpha(z)=\mathbf 1_{\{z_i=t\}}\) and
\(\beta(u)=g_i(t,u)\); the second equality holds because the indicator
vanishes off \(\{U_{0,i}=t\}\), and on that event the first argument of
\(g_i\) is \(t\). The randomization property and (F3), with \(U_i\) a
coordinate of \(\bar U\), make \(U_0\) and \(U_i\) independent. By (F2),
the expectation equals
\(\mathbb E[\alpha(U_0)]\,\mathbb E[\beta(U_i)]
 =\pi_i(t)\,\mathbb E[g_i(t,U_i)]\), since
\(\mathbb E[\alpha(U_0)]=\sum_{z:z_i=t}\varpi(z)=\pi_i(t)\).
Dividing by \(\pi_i(t)>0\) and averaging over \(i\) gives
\(\Phi_t(p_M)=\mu_M(t)\). Multiplying by \(w_t\) and summing gives
\(\tau(M)=\Phi_w(p_M)\).

The probabilities \(\pi_i(t)\) are fixed by \(\varpi\), hence common to
the whole class. If \(M_1,M_2\in\mathcal C_{\mathrm{NI}}(\varpi)\) have
\(p_{M_1}=p_{M_2}\), then
\(\tau(M_1)=\Phi_w(p_{M_1})=\Phi_w(p_{M_2})=\tau(M_2)\). This is
identification in the sense of \cref{def:p16-causal-identification}, with
all vertices observed and \(Q=\tau\).

\emph{Part 2.} For \(v=(z,y)\), only the term at \((z,y)\) survives in
\(\Phi_w(\delta_v)\), giving
\(\sum_tw_tN^{-1}\sum_i\mathbf 1_{\{z_i=t\}}y_i/\pi_i(t)\). At
\(v=V(\omega)\) this is \(\widehat\tau(\omega)\). The map \(\Phi_w\) is a
finite sum of values of \(p\) with fixed coefficients, so
\(\Phi_w(\sum_kc_kp_k)=\sum_kc_k\Phi_w(p_k)\) for finitely many finitely
supported \(p_k\). Evaluating at any \(v\) shows
\(p_M=\sum_\omega q(\omega)\delta_{V(\omega)}\). Therefore
\[
 \mathbb E[\widehat\tau]=\sum_\omega q(\omega)\Phi_w(\delta_{V(\omega)})
 =\Phi_w\Bigl(\sum_\omega q(\omega)\delta_{V(\omega)}\Bigr)
 =\Phi_w(p_M)=\tau(M).
\]
For \(M=M_E\), Part~1 of \cref{thm:p15-w02b-p15-077-correspondence} gives
\(V_i=Y_i(Z_i)\), so \(\widehat\tau\) is the estimator of
\cref{thm:p15-factorial-assignment}, and this conclusion agrees with the
unbiasedness proved there.

\emph{Part 3.} Let \(\mathcal Z\) be the set of vectors with exactly
\(n_s\) coordinates equal to \(s\), for every \(s\), so that \(\varpi\)
is uniform on \(\mathcal Z\). For \(i\ne j\), exchanging coordinates
\(i\) and \(j\) maps \(\mathcal Z\) bijectively onto itself and maps
\(\{z\in\mathcal Z:z_i=t\}\) onto \(\{z\in\mathcal Z:z_j=t\}\). Under the
uniform law these sets have equal probability, so \(\pi_i(t)\) does not
depend on \(i\). Every \(z\in\mathcal Z\) satisfies
\(\sum_i\mathbf 1_{\{z_i=t\}}=n_t\), so
\(\sum_i\pi_i(t)=\mathbb E[\sum_i\mathbf 1_{\{Z_i=t\}}]=n_t\).
Hence \(\pi_i(t)=n_t/N>0\). Since \(n_t\) is fixed,
\(\overline Y_t=n_t^{-1}\sum_i\mathbf 1_{\{Z_i=t\}}V_i\) and, by the
computation in Part~1,
\[
 \mathbb E[\overline Y_t]=\frac1{n_t}\sum_i\pi_i(t)\,
   \mathbb E[g_i(t,U_i)]
 =\frac1N\sum_i\mathbb E[g_i(t,U_i)]=\mu_M(t).
\]
Multiplying by \(w_t\) and summing gives the formula for \(\tau(M)\).
\end{proof}
\end{annexitem}

\begin{annexitem}{annex:p15-w02b-p15-077-interf}{thm:p15-w02b-p15-077-interference}{Interference}
\begin{proof}
Let \(M\in\mathcal C_{\mathrm{int}}(\varpi)\). Its graph and order are
those of the preceding proofs, and observationally \(V_0=U_0\),
\(V_i=f_i(U_0,U_i)\).

\emph{Part 1.} Under \(\operatorname{do}(V_0=t\mathbf 1)\),
\(V_i^{0=t\mathbf 1}=f_i(t\mathbf 1,U_i)\), so
\(\mu_M(t)=\mathbb E[H(\bar U)]\) with
\(H(u)=N^{-1}\sum_if_i(t\mathbf 1,u_i)\). Observationally, on the event
\(\{U_0=t\mathbf 1\}\) each \(V_i\) equals \(f_i(t\mathbf 1,U_i)\),
so
\[
 \mathbb E\bigl[h(V_1,\ldots,V_N)\mathbf 1_{\{V_0=t\mathbf 1\}}\bigr]
 =\mathbb E\bigl[\mathbf 1_{\{U_0=t\mathbf 1\}}H(\bar U)\bigr]
 =\varpi(t\mathbf 1)\,\mathbb E[H(\bar U)],
\]
by the randomization property and (F2). Since
\(\mathbb P(V_0=t\mathbf 1)=\varpi(t\mathbf 1)>0\), dividing gives the
stated conditional expectation, which is a functional of \(p_M\) by (F1).

\emph{Part 2.} Define \(M'\) from \(M\) by changing only vertex \(1\):
\(S_1'=S_1\cup\{s+c:s\in S_1\}\) and
\(f_1'(z,u)=f_1(z,u)+c\,\mathbf 1_{\{z=t^*\mathbf 1\}}\). The function
is total on its finite domain, and properties 1 and 2 of
\cref{def:p15-w02b-p15-077-classes} involve only \(f_0\), \(U_0\) and
\(\bar U\), which are unchanged; thus
\(M'\in\mathcal C_{\mathrm{int}}(\varpi)\). Let
\(A=\{\omega:U_0(\omega)=t^*\mathbf 1\}\). Then
\(\sum_{\omega\in A}q(\omega)=\varpi(t^*\mathbf 1)=0\), and since every
\(q(\omega)\geq0\), \(q\) vanishes on \(A\). For \(\omega\notin A\),
\(V_0(\omega)\ne t^*\mathbf 1\), so evaluation in the topological order
gives \(V'(\omega)=V(\omega)\). Hence, for every \(v\),
\[
 p_{M'}(v)=\sum_{\omega\notin A:\,V(\omega)=v}q(\omega)
          +\sum_{\omega\in A:\,V'(\omega)=v}q(\omega)
          =\sum_{\omega:\,V(\omega)=v}q(\omega)=p_M(v),
\]
because both sums over \(A\) are zero. Under
\(\operatorname{do}(V_0=t^*\mathbf 1)\), \(V_1'=V_1+c\) at every outcome
and \(V_i'=V_i\) for \(i\geq2\), so \(\mu_{M'}(t^*)=\mu_M(t^*)+c/N\).
For \(t\ne t^*\), the vectors \(t\mathbf 1\) and \(t^*\mathbf 1\) differ
in their first coordinate, so \(f_1'(t\mathbf 1,\cdot)=f_1(t\mathbf 1,\cdot)\)
and \(\mu_{M'}(t)=\mu_M(t)\).

\emph{Part 3.} If \(\varpi(t\mathbf 1)>0\) whenever \(w_t\ne0\), Part~1
expresses \(\tau(M)=\sum_{t:w_t\ne0}w_t
\mathbb E[h(V_1,\ldots,V_N)\mid V_0=t\mathbf 1]\) through \(p_M\) alone,
so models with equal observational laws have equal \(\tau\). Conversely,
let \(w_{t^*}\ne0\) and \(\varpi(t^*\mathbf 1)=0\). The class is
nonempty: the structural model of the experiment with every
\(Y_i(t)=0\) and design law \(\varpi\) belongs to
\(\mathcal C_{\mathrm{NI}}(\varpi)\subseteq\mathcal C_{\mathrm{int}}(\varpi)\).
For any \(M\) in the class, Part~2 with \(c=1\) gives \(M'\) with
\(p_{M'}=p_M\) and \(\tau(M')-\tau(M)=w_{t^*}/N\ne0\). By
\cref{def:p16-causal-identification}, \(\tau\) is not identified.

Under complete randomization with at least two cells of positive size,
every vector in the support of \(\varpi\) has coordinates taking at least
two values, so \(t\mathbf 1\) lies outside the support and
\(\varpi(t\mathbf 1)=0\) for every \(t\). Any nonzero weight vector has
some \(w_{t^*}\ne0\), so \(\tau\) is not identified. In the pair
\(M,M'\) just constructed, \(\Phi_w(p_M)=\Phi_w(p_{M'})\) while
\(\tau(M)\ne\tau(M')\); hence \(\Phi_w(p)\ne\tau\) for at least one of
the two models.
\end{proof}
\end{annexitem}

\begin{annexitem}{annex:p15-w02b-p15-077-noncomp}{thm:p15-w02b-p15-077-noncompliance}{Noncompliance}
\begin{proof}
Every edge goes from \(0\) to some \(d_i\) or from \(d_i\) to
\(y_i\), so the graph is acyclic, with topological order \(0\), then all
\(d_i\), then all \(y_i\).

\emph{Part 1.} Put \(W_i=(U_{d_i},U_{y_i})\) and
\(G_i(t,w)=f_{y_i}(c_i(t,w_1),w_2)\). Under
\(\operatorname{do}(V_0=t\mathbf 1)\), only vertex \(0\) changes, so
\(Y_i^{0=t\mathbf 1}=G_i(t,W_i)\) and
\(\mu_M(t)=N^{-1}\sum_i\mathbb E[G_i(t,W_i)]\). Observationally,
\(Z=U_0\) and \(Y_i=G_i(Z_i,W_i)\). The tuple \(W_i\) is a function of
the tuple of all disturbances other than \(U_0\), so by (F3) it is
independent of \(U_0\). The argument proving Part~1 of
\cref{thm:p15-w02b-p15-077-identification}, with \(g_i\) replaced by
\(G_i\) and \(U_i\) by \(W_i\), gives
\(\mathbb E[\mathbf 1_{\{Z_i=t\}}Y_i]=\pi_i(t)\mathbb E[G_i(t,W_i)]\).
The left side depends only on the marginal law of
\((V_0,Y_1,\ldots,Y_N)\), by (F1). Dividing by \(\pi_i(t)\) and averaging
gives \(\mu_M(t)=\Phi_t(p_{M,\{0,y_1,\ldots,y_N\}})\), so the assignment
contrast is identified within \(\mathcal K(\varpi)\).

\emph{Part 2.} Take \(\Omega=\{0,1\}^N\times\{0,1\}^N\), with
outcomes \((z,v)\) and mass \(q(z,v)=\varpi(z)2^{-N}\); these masses are
nonnegative and sum to one. Put \(U_0(z,v)=z\) and
\(U_{d_i}(z,v)=U_{y_i}(z,v)=v_i\). Then
\(\mathbb P(U_0=z)=\sum_vq(z,v)=\varpi(z)\), and
\(\mathbb P(U_0=z,(v_i)_i=v')=\varpi(z)2^{-N}\) shows that \(U_0\) is
independent of the tuple \((v_i)_i\), hence by (F3) of the tuple of all
other disturbances. Let \(c_i(t,u)=tu\). In \(M_A\) let
\(f_{y_i}(\delta,u)=u\); in \(M_B\) let
\(f_{y_i}(\delta,u)=\max(\delta,u)\). Both have
\(S_{y_i}=\{0,1\}\), and both lie in \(\mathcal K(\varpi)\): the uptake
function reads only \(Z_i\) and the response function reads only
\(D_i\) besides its disturbance. The disturbances of \(d_i\) and \(y_i\)
coincide; the class permits this dependence.

At \(\omega=(z,v)\), both models give \(Z=z\) and \(D_i=z_iv_i\). In
\(M_A\), \(Y_i=v_i\). In \(M_B\),
\(Y_i=\max(z_iv_i,v_i)=v_i\), because \(z_iv_i\leq v_i\). The two
solutions coincide at every outcome, so the observational laws on all
vertices are equal. Further,
\(\mathbb P(D_i=1)=\mathbb P(Z_i=1,v_i=1)=\pi_i(1)/2>0\) and
\(\mathbb P(D_i=0)\geq\mathbb P(Z_i=0)=\pi_i(0)>0\).

Under \(\operatorname{do}(V_0=t\mathbf 1)\), \(D_i=tv_i\); then
\(Y_i=v_i\) in \(M_A\) and \(Y_i=\max(tv_i,v_i)=v_i\) in \(M_B\). Since
\(\mathbb E[v_i]=\sum_{z,v}\varpi(z)2^{-N}v_i=1/2\), both models have
\(\mu(0)=\mu(1)=1/2\) and assignment contrast zero. Under the intervention
setting every \(d_i\) to \(s\), \(Y_i=v_i\) in \(M_A\), so
\(\nu_{M_A}(1)=\nu_{M_A}(0)=1/2\). In \(M_B\), \(Y_i=\max(s,v_i)\),
which is \(1\) for \(s=1\) and \(v_i\) for \(s=0\); hence
\(\nu_{M_B}(1)-\nu_{M_B}(0)=1-1/2=1/2\). Two models of the class with
equal observational laws and unequal received-treatment contrasts prove
nonidentification by \cref{def:p16-causal-identification}.
\end{proof}
\end{annexitem}

\begin{annexitem}{annex:p15-w02b-p15-077-miss}{thm:p15-w02b-p15-077-missingness}{Missing outcomes}
\begin{proof}
Every edge goes from \(0\) to some \(y_i\) or \(r_i\), from \(y_i\)
to \(r_i\) or \(o_i\), or from \(r_i\) to \(o_i\); the order \(0\),
all \(y_i\), all \(r_i\), all \(o_i\) is topological, so the graph is
acyclic.

\emph{Part 1.} Use \(\Omega\), \(q\) and \(U_0\) of the preceding proof,
and write \(s_i=v_i\). Value sets are \(S_{y_i}=\{0,1\}\),
\(S_{r_i}=\{0,1\}\) and \(S_{o_i}=\{0,1,\star\}\).
In \(M_A\), let \(U_{y_i}\) be constant, \(g_i(t,u)=t\),
\(U_{r_i}=s_i\), and \(\rho_i(t,y,u)=1\) for \(t=0\) and \(u\) for
\(t=1\). In \(M_B\), let \(U_{y_i}=s_i\), \(g_i(t,u)=tu\),
\(U_{r_i}\) be constant, and \(\rho_i(t,y,u)=1\) for \(t=0\) and \(y\)
for \(t=1\). In both models the recording indicator takes values in
\(\{0,1\}\), and \(U_0\) is independent of the tuple of other
disturbances, which are functions of \((s_i)_i\), by (F3) as before.

At \(\omega=(z,s)\), a unit with \(z_i=0\) has \(Y_i=0\), \(R_i=1\) and
\(O_i=0\) in both models. A unit with \(z_i=1\) has, in \(M_A\),
\(Y_i=1\), \(R_i=s_i\); in \(M_B\), \(Y_i=s_i\) and
\(R_i=Y_i=s_i\). In both, \(O_i=1\) when \(s_i=1\) and \(O_i=\star\)
when \(s_i=0\). The observed coordinates therefore coincide at every
outcome, and so do the observed laws
\(p_{M,O}(o)=\sum_{\omega:V_O(\omega)=o}q(\omega)\); in particular the
arm-specific recording probabilities agree. Under
\(\operatorname{do}(V_0=t\mathbf 1)\), \(M_A\) gives \(Y_i=t\) and
\(M_B\) gives \(Y_i=ts_i\), with \(\mathbb E[s_i]=1/2\). Hence
\(\tau(M_A)=1-0=1\) and \(\tau(M_B)=1/2-0=1/2\).

\(M_A\) lies in \(\mathcal W_0(\varpi)\): its \(\rho_i\) does not read
\(y\); for each \(i\),
\(\mathbb P(U_0=z,U_{y_i}=0,s_i=a)=\sum_{v:v_i=a}\varpi(z)2^{-N}
=\varpi(z)/2\), the product of the three marginals; and
\(\bar\rho_i(0)=1\), \(\bar\rho_i(1)=\mathbb P(s_i=1)=1/2\).

\emph{Part 2.} Let \(M\in\mathcal W_0(\varpi)\), and fix \(i\) and \(t\).
Under \(\operatorname{do}(V_0=t\mathbf 1)\) the equation of \(y_i\) is
unchanged, so \(Y_i^{0=t\mathbf 1}=g_i(t,U_{y_i})\) and
\(\mu_M(t)=N^{-1}\sum_i\mathbb E[g_i(t,U_{y_i})]\). Observationally
\(Z=U_0\), \(Y_i=g_i(Z_i,U_{y_i})\) and
\(R_i=\rho_i(Z_i,U_{r_i})\). When \(R_i=1\), \(O_i=Y_i\); when
\(R_i=0\), \(R_iO_i\) is read as zero. With
\(\alpha(z)=\mathbf 1_{\{z_i=t\}}\), and arguing on the event
\(\{Z_i=t\}\) as in the identification proof,
\[
 \mathbf 1_{\{Z_i=t\}}R_iO_i
   =\alpha(U_0)\,\rho_i(t,U_{r_i})\,g_i(t,U_{y_i}),\qquad
 \mathbf 1_{\{Z_i=t,R_i=1\}}=\alpha(U_0)\,\rho_i(t,U_{r_i}).
\]
By mutual independence of the three disturbances
\(U_0\), \(U_{y_i}\) and \(U_{r_i}\), and by (F2), the two expectations
are
\[
 \pi_i(t)\,\bar\rho_i(t)\,\mathbb E[g_i(t,U_{y_i})]
 \quad\text{and}\quad
 \pi_i(t)\,\bar\rho_i(t),
\]
using
\(\mathbb E[\rho_i(t,U_{r_i})]=\bar\rho_i(t)\) for a
\(\{0,1\}\)-valued function. The second is positive. Their ratio is
\(\mathbb E[g_i(t,U_{y_i})]\); averaging over \(i\) gives the displayed
formula. Both expectations are functionals of the law of the observed
coordinates \(Z,R_i,O_i\), by (F1), so models of
\(\mathcal W_0(\varpi)\) with equal observed laws have equal
\(\mu_M(t)\) for every \(t\), and equal \(\tau\).
\end{proof}
\end{annexitem}

\begin{annexitem}{annex:p15-w02b-p15-077-study}{thm:p15-w02b-p15-077-study-contrasts}{Declared contrasts}
\begin{proof}
Identify condition \((a,b)\) with the cell \(t=ab\) of
\cref{def:p15-factorial-coding}, in the order \((00,10,01,11)\). The
aid-present weights \((0,0,-1,1)\) give
\(\tau=\mu(11)-\mu(01)\); the equal-weight weights
\((-1,1,-1,1)/2\) give
\(\tau=\bigl((\mu(10)-\mu(00))+(\mu(11)-\mu(01))\bigr)/2\). Both are
identified within \(\mathcal C_{\mathrm{NI}}(\varpi)\) by Part~1 of
\cref{thm:p15-w02b-p15-077-identification}. Complete randomization with
fixed positive cell sizes is the design of its Part~3, which gives
\(\mu(t)=\mathbb E[\overline Y_t]\) for each cell. Linearity of
expectation then gives the two displayed expectations. In this class every
response is recorded, so the observed condition means
\(\overline Y_{ab}\) of \cref{def:p19-proposed-factorial-design} are the
assigned-arm means \(\overline Y_t\). Hence
\(\widehat\Delta_{\rm aid}=\overline Y_{11}-\overline Y_{01}\) and
\(\widehat\Delta_{\rm check}\) have expectations equal to the two
contrasts, which is their unbiasedness over the model's probability law,
in particular over the randomization when the responses are fixed.
\end{proof}
\end{annexitem}
